\section{Interpolation with a common small parameter}
\label{int:section}

The height theorem below applies to equations assembled from a fixed
collection of ordinary differential systems.  Its hypotheses distinguish
the global height of the target from the much smaller positive coordinate
logarithms at the places used for interpolation.  We state the theorem
and explain the graph construction that exploits this distinction here.
The full proof is given in Appendix~\ref{int:proof}.

\subsection{The quantitative statement}

For a tuple $x=(x_1,\ldots,x_r)$ over a number field $F$, write
\[
 h_{\mathrm{aff}}(x)
   =\sum_v w_v\log\max(1,|x_1|_v,\ldots,|x_r|_v),
 \qquad w_v=\frac{[F_v:\Q_v]}{[F:\Q]}.
\]
At an archimedean place we use the ordinary absolute value, so that a complex
place has weight $2/[F:\Q]$. At a finite place above $p$, the absolute
value extends $|p|_p=p^{-1}$. The same notation for a polynomial means the
height of its coefficient tuple.  A weighted subcollection of places may
split a place into its embeddings, with the corresponding weights.  When
$F$ is enlarged, these weights are split among the extensions of each
embedding.  In particular all sums for quantities defined over the old
field are unchanged.  We write $\sum_V w_v$ with this convention even when
the subcollection is described in terms of embeddings.

An \emph{ordinary germ} below is a horizontal analytic solution in a rational
frame in which the connection is regular at the specified center.  This is
a condition on a differential equation.  The rank, frame, center, and curve
carrying each such equation are fixed.

\begin{theorem}[Conditional height bound]\label{int:height}
Fix an ambient dimension $N$, a bound $s_0$ for the number of equations,
an integer $b$ bounding the degrees of the polynomial combinations below,
and a real constant $C_0\geq2$.  Fix also a finite list of algebraic regular
germs and ordinary horizontal matrix germs on fixed algebraic curves over
number fields.  Dual and inverse matrices are allowed in this list.
Then there is a constant $C$, depending only on these fixed data, with the
following property.

Let $d,h\geq2$ and let $z=(a,U)\in\mathbb A^N(\Qbar)$ with $a\ne0$.
Suppose a weighted subcollection $V$ and $1\leq s\leq s_0$ equations
\begin{equation}\label{int:equations}
 E_\alpha(a,U)=\sum_{j\geq0}a^j e_{\alpha,j}(U)=0
 \quad(1\leq\alpha\leq s)
\end{equation}
satisfy the following conditions.
\begin{enumerate}[label=\textup{(\roman*)}]
\item The point and the selected places satisfy
\begin{align}
 h_{\mathrm{aff}}(z)&\leq C_0d^{C_0}h,
 &|a|_v&<1\quad(v\in V),\label{int:global-input}\\
 \sum_{v\in V}w_v\log |a|_v^{-1}
   &\geq \frac{h}{C_0d^{C_0}},
 &\sum_{v\in V}w_v\log^+\|U\|_v
   &\leq C_0d^{C_0}(1+\log h)^{C_0}.
 \label{int:local-input}
\end{align}
Here $\|U\|_v$ is the maximum of the coordinate absolute values.
\item The coefficients are polynomials over a number field containing the
target and equation data, with
\begin{equation}\label{int:coefficient-input}
 \deg e_{\alpha,j}\leq C_0(1+j),\qquad
 h_{\mathrm{aff}}(e_{\alpha,j})
  \leq C_0(1+j+d)^{C_0}(1+\log h)^{C_0}.
\end{equation}
The same equations are used at every selected embedding.  They converge
there at the target, and there are numbers $b_v\geq0$ such that, for all
$\ell\geq0$ and all $\alpha$,
\begin{align}
 \left|\sum_{j\geq\ell}a^j e_{\alpha,j}(U)\right|_v
   &\leq |a|_v^\ell\exp((\ell+1)b_v),\label{int:tail-input}\\
 \sum_{v\in V}w_vb_v
   &\leq C_0d^{C_0}(1+\log h)^{C_0}.
\end{align}
\item As analytic functions in the variables $(a,U)$, the $E_\alpha$ are
polynomial combinations of degree at most $b$ of the coordinates and the
fixed germs evaluated at $a$, and possibly at $aR$, where $R$ is one of
the coordinates of $U$.  Their constant coefficients may depend on the
target.  The algebraic germs are regular functions on the specified local
curve charts.  Thus every expression is also a convergent complex analytic
germ near any finite point of $a=0$.
\item Let $W\subseteq\mathbb A^N_{\Qbar}$ be any irreducible closed
subvariety containing $z$ on which $a$ is nonconstant.  For each prime
divisor $D$ in its finite normalization lying above a component of
$W\cap\{a=0\}$, at least one $E_\alpha$ is not identically zero as an
analytic germ on the local branch at a general point of $D$.
\end{enumerate}
Then $h\leq Cd^C$.
\end{theorem}

The quantifier in condition \textup{(iv)} includes varieties defined using
the coefficients of the equations and the auxiliary polynomials constructed
below.  It concerns the germ on a whole local branch, rather than its value
on the boundary divisor.  No bound on the degree of the field of equation
coefficients is part of this theorem.

For the elliptic-square application, the inputs to conditions
\textup{(i)--(iii)} are assembled in \Cref{loc:groups}, and condition
\textup{(iv)} is verified in \Cref{nid:nonidentity}; the resulting height
bound is stated in \Cref{nid:height}.  We next explain why the theorem
requires nonidentity on every algebraic branch, and why small positive
local logarithms are needed in addition to a global height bound.

\subsection{Why the graph gives a height bound}

The proof constructs an auxiliary polynomial of degree $L$ whose value
at the target gains a factor $a^n$ at the selected places.  To make the
product formula useful, the construction must allow a sufficiently
large ratio $n/L$.
The graph construction makes this possible by lowering the dimension
of the boundary on which the interpolation conditions are imposed.

Start with $W=\mathbb A^N$.  During the proof, $W$ will be an irreducible
algebraic subvariety containing $z$, with controlled degree and defining
heights.  Write $k=\dim W$ and first suppose that $a$ varies on $W$ and
that $W\cap\{a=0\}$ is nonempty.  Condition \textup{(iv)} and an ordinary
differential multiplicity estimate give an integral linear combination
$E=\sum c_\alpha E_\alpha=\sum_{j\geq0}a^je_j(U)$ and an integer $m$
such that its truncation
\[
 F=E_{<m}:=\sum_{j=0}^{m-1}a^je_j(U)
 \quad\text{satisfies}\quad
 \ord_D F<m\ord_D a
\]
on every normalized boundary divisor $D$.  The coefficients $c_\alpha$
and $m$ are polynomially bounded in $\deg W$.  This strict inequality
means that $F/a^m$ has a pole at the generic point of every such divisor;
a nonzero restriction of $E$ to the boundary itself is not required.

Let $W'$ be the affine closure of the graph of $q=F/a^m$ over $a\ne0$.
It has dimension $k$.  Its boundary over $a=0$ has no point above a general
point of any boundary divisor of $W$, because $q$ has a pole there.
Thus the projection of $W'\cap\{a=0\}$ to $W$ has dimension at most
$k-2$.  Imposing $q=-e_m(U)$ fixes the graph coordinate, so
\[
 \dim\bigl(W'\cap\{a=0,\ q+e_m(U)=0\}\bigr)\leq k-2.
\]
For $k=1$ this intersection is empty.  This is a bound on the indicated
incidence, not on the whole graph boundary, whose fibers may have
positive dimension.

The lifted target $z'=(a,U,F(z)/a^m)$ has a different useful property.
At each selected place $E(z)=0$, and therefore
\[
 q(z')=-a^{-m}\sum_{j\geq m}a^je_j(U).
\]
The tail estimate bounds the positive local logarithm of this new
coordinate by $(m+1)b_v$, with the errors enlarged to accommodate the
chosen integral combination.  Hence $z'$ still has small positive
local logarithms, although its global height can be proportional to $h$.

Choose equations generating an ideal $I$ whose zero set is $W'$ and put
\[
 F_n=q+\sum_{j=0}^{n-1}a^je_{m+j}(U).
\]
The dimension bound above permits a polynomial $P$, nonzero on $W'$,
with
\[
 \deg P\leq L,\qquad P\in(I,a^n,F_n),\qquad n=cL,
\]
for any prescribed positive integer $c$, with $L$ polynomially bounded
in $c$ and the degree bounds of the construction.  Here membership is
in the displayed ideal itself, even if $I$ is nonradical.  The reason
for the freedom in $c$ is a count of linear conditions.  For $k\geq2$,
reducing successively through $n$ powers of $a$ imposes
$O(n(L+n)^{k-2})$ conditions, with a constant polynomial in the defining
degree bounds.  Restrictions of degree-$\leq L$ polynomials to $W'$
span a space of dimension at least $\binom{L+k}{k}$.  For $n=cL$ these
grow respectively as $L^{k-1}$ and $L^k$.  When $k=1$ there are no
conditions.  This saving of one power of $L$ is the purpose of the graph.

At the selected places, evaluating the ideal-membership certificate and
using the tail bound gives the factor $|a|_v^n$.  Even if the certificate
has degree larger than $L$, its degree multiplies only the small
positive local logarithms of $z'$.  At the remaining places one evaluates
$P$ itself, so the global height cost is proportional to $Lh$.
Consequently, if $P(z')\ne0$, the product formula gives an inequality
of the form
\[
 0\leq\left(-\frac n{A(d)}+L C(d)\right)h
        +D(d)(1+\log h)^B,
\]
where $A,C,D$ are fixed polynomials and $B$ is fixed.  The polynomials
$A,C$ are determined before $c$ is chosen.  Taking $c$ sufficiently
large, still polynomially bounded in $d$, makes the coefficient of $h$
negative and yields a polynomial bound for $h$ in $d$.

If that bound does not hold, $P(z')=0$.  Substituting $q=F/a^m$ and
clearing denominators gives a polynomial nonzero on $W$ and zero at
$z$.  An irreducible component through $z$ of its zero set in $W$
has smaller dimension and controlled degree and height.  This explains
the universal quantifier in condition \textup{(iv)}: the next variety
is chosen using the equations and the auxiliary polynomial.  After at
most $N$ such steps, either the height is bounded or one reaches a
variety disjoint from $a=0$.  If $G_1,\ldots,G_r$ define that variety,
a polynomial identity $1=aA+\sum_\rho G_\rho A_\rho$ and the small
positive local logarithms give the bound directly.
Appendix~\ref{int:proof} proves these degree, height, and multiplicity
assertions and the descent in detail.
